\documentclass[../sablona-main.tex]{subfiles}

\begin{document}
	\section{Grafické řešení pro změněný zisk (Úkol 8)}
	Po vykreslení množiny splňující všechny vazebné podmínky a nové účelové funkce $120x + 300y = z$ obdržíme maximální hodnotu zisku pro $z = 12600$. Celá situace je znázorněna na obrázku \ref{fig: Graficke reseni alternativni}.
	
	\begin{figure}[H]
		\centering
		\begin{tikzpicture}[line cap=round,line join=round,>=triangle 45]
	\begin{axis}[
		x=1.5mm, y=1.5mm,
		axis lines=middle,
		ymajorgrids=true,
		xmajorgrids=true,
		grid style={line width=.1pt, draw=gray!30},
		xmin=-18, xmax=75,
		ymin=-11, ymax=43,
		xtick={-15,-10,...,75},
		ytick={-10,-5,...,40},
		]
		% Zelená plocha - omezující podmínky (zůstává beze změny)
		\fill[color=qqffqq, fill opacity=0.3] 
		(0,35) -- (0,0) -- (50,0) -- (50,10) -- (30,30) -- (17.5,35) -- cycle;
		
		% Obrys zelené plochy
		\draw[line width=1.5pt, color=qqffqq] 
		(0,35) -- (0,0) -- (50,0) -- (50,10) -- (30,30) -- (17.5,35) -- cycle;
		
		% Červená čára - nová účelová funkce pro Z = 12600
		\draw [line width=2pt, color=ffqqqq, domain=-18:75, samples=2] 
		plot(\x,{(12600-120*\x)/300});

		% Zvýraznění úsečky alternativních optimálních řešení (mezi body A a B)
		\draw [line width=3pt, color=blue, dashed] 
		(17.5,35) -- (30,30);
		
	\end{axis}
\end{tikzpicture}
		\caption{Grafické řešení s alternativními optimálními řešeními} \label{fig: Graficke reseni alternativni}
	\end{figure}
	
	Z obrázku je zřejmé, že zisková přímka je rovnoběžná s hranicí omezující podmínky pro montážní hodiny. Tento jev nastává proto, že obě lineární rovnice mají totožnou směrnici. 
	
	Rovnice účelové funkce:
	$$120x + 300y = z \implies y = -0.4x + \frac{z}{300}$$
	
	a pro omezující podmínku montážních hodin dostaneme:
	$$0.6x + 1.5y = 63 \implies y = -0.4x + 42$$
	
	Globálním maximem tak není pouze jeden bod (vrchol), ale všechny body ležící na úsečce, která spojuje vrcholy $(17.5, 35)$ a $(30, 30)$. Každá kombinace výroby na této úsečce přináší naprosto stejný, maximální možný zisk.
\end{document}